\documentclass[14pt,a4paper]{extarticle}

% Переносимо: компилируется и pdflatex (T2A), и xelatex/lualatex (fontspec).
\usepackage{iftex}
\ifPDFTeX
    \usepackage[T2A]{fontenc}
    \usepackage[utf8]{inputenc}
\else
    \usepackage{fontspec}
    \setmainfont{Liberation Serif}
\fi
\usepackage[russian]{babel}
\usepackage[left=30mm,right=15mm,top=20mm,bottom=20mm]{geometry}
\usepackage{graphicx}
\usepackage{hyperref}
\usepackage{listings}
\usepackage{xcolor}

\newcommand{\group}{Р3413}
\newcommand{\worknum}{1}
\newcommand{\workname}{Разработка защищённого REST API с интеграцией в CI/CD}

\lstset{
    basicstyle=\ttfamily\small,
    breaklines=true,
    frame=single,
    columns=fullflexible,
    keepspaces=true,
    showstringspaces=false
}

\newcommand{\shot}[1]{%
    \begin{center}%
    \fbox{\includegraphics[width=0.98\linewidth,height=0.42\textheight,keepaspectratio]{#1}}%
    \end{center}%
}

\begin{document}

\begin{titlepage}
    \begin{center}
        Федеральное государственное автономное образовательное учреждение\\
        высшего образования\\
        «Национальный исследовательский университет ИТМО»

        \vfill

        Лабораторная работа №\worknum

        \vspace{0.5cm}
        по дисциплине «Информационная безопасность»

        \vspace{0.5cm}
        \textbf{\workname}
    \end{center}

    \vfill

    \begin{flushright}
        Выполнил:\\
        Карпов Александр Дмитриевич\\
        \group
    \end{flushright}

    \vfill
    \begin{center}
        Санкт-Петербург, 2026
    \end{center}
\end{titlepage}

\tableofcontents
\newpage

\section{Цель работы}
Получить практический опыт разработки безопасного backend-приложения с
автоматизированной проверкой кода на уязвимости. Освоить базовые меры защиты из
OWASP Top 10 (инъекции, XSS, недостатки аутентификации) и интеграцию
инструментов анализа безопасности в CI/CD.

\section{Постановка задачи}
Разработать REST API со следующими эндпоинтами:
\begin{itemize}
    \item \texttt{POST /auth/login} — аутентификация пользователя;
    \item \texttt{GET /api/data} — получение данных, доступно только
          аутентифицированным пользователям;
    \item \texttt{POST /api/buy} — собственный метод: покупка кристаллов
          (дополнительно реализован \texttt{POST /auth/register}).
\end{itemize}
Реализовать защиту от SQL-инъекций, XSS и небезопасной аутентификации; настроить
CI/CD pipeline с запуском SAST- и SCA-сканеров.

\section{Используемые инструменты}
\begin{itemize}
    \item Python 3.12, Flask, Flask-SQLAlchemy, SQLite;
    \item PyJWT — выпуск и проверка JWT;
    \item bcrypt — хэширование паролей;
    \item Bandit — статический анализ (SAST);
    \item pip-audit — анализ зависимостей (SCA);
    \item GitHub Actions — CI/CD;
    \item pytest — автотесты; curl — ручное тестирование.
\end{itemize}

\section{Ход выполнения}

\subsection{Архитектура приложения}
Приложение собрано по схеме application factory. Логика вынесена в пакет
\texttt{app}: модели (\texttt{models.py}), выпуск и проверка токена
(\texttt{auth.py}), маршруты (\texttt{routes.py}), начальное наполнение БД
(\texttt{seed.py}). Точка входа \texttt{run.py} лишь создаёт приложение.

\subsection{Защита от SQL-инъекций}
Все обращения к БД выполняются через ORM SQLAlchemy
(\texttt{filter\_by}, \texttt{db.session.get}). Пользовательский ввод всегда
передаётся как параметр запроса и экранируется драйвером, а не склеивается в
строку SQL. Попытка входа с пейлоадом \texttt{admin' OR '1'='1} приводит к ответу
401 — инъекция не срабатывает.

\subsection{Защита от XSS}
Все строковые поля, пришедшие от пользователя (\texttt{username}, \texttt{color}),
экранируются функцией \texttt{markupsafe.escape} перед возвратом в JSON. Имя
вида \texttt{<script>alert(1)</script>} возвращается как
\texttt{\&lt;script\&gt;...} и не исполняется на клиенте.

\subsection{Аутентификация}
Пароли хранятся только в виде bcrypt-хэша со случайной солью
(\texttt{bcrypt.hashpw}). При успешном входе выдаётся подписанный JWT (алгоритм
HS256) с временем жизни. Защищённые эндпоинты закрыты декоратором
\texttt{token\_required}, который проверяет наличие заголовка
\texttt{Authorization: Bearer}, подпись токена и срок действия (\texttt{exp}).

\subsection{CI/CD}
Файл \texttt{.github/workflows/ci.yml} запускается на каждый \texttt{push} и
\texttt{pull\_request} и последовательно выполняет: установку зависимостей,
Bandit (SAST), pip-audit (SCA) и pytest.

\section{Результаты}

\subsection{Работа API}

\paragraph{Логин и получение токена.}
\shot{images/01_login.png}

\paragraph{Доступ к \texttt{/api/data} без токена (401) и с токеном (200).}
\shot{images/02_data_auth.png}

\paragraph{SQL-инъекция в логине отклонена.}
\shot{images/03_sqli.png}

\paragraph{XSS: имя пользователя экранировано в ответе.}
\shot{images/04_xss.png}

\paragraph{Покупка кристаллов.}
\shot{images/05_buy.png}

\subsection{Отчёты SAST/SCA}

\paragraph{Bandit (SAST) — критических проблем не найдено.}
\shot{images/06_bandit.png}

\paragraph{pip-audit (SCA) — уязвимых зависимостей не найдено.}
\shot{images/07_pipaudit.png}

\paragraph{Автотесты (pytest) в составе pipeline.}
\shot{images/09_tests.png}

\paragraph{Успешный запуск pipeline (все шаги зелёные).}
\shot{images/08_actions.png}

\section{Ссылки}
\begin{itemize}
    \item Репозиторий: \url{https://github.com/Alexander-D-Karpov/ib-lab1}
    \item Последний успешный pipeline: \\
          \url{https://github.com/Alexander-D-Karpov/ib-lab1/actions/runs/37653150323}
\end{itemize}

\section{Вывод}
Разработан защищённый REST API с аутентификацией по JWT и хэшированием паролей
bcrypt. Реализована защита от SQL-инъекций (ORM с параметризацией) и XSS
(экранирование вывода). Настроен CI/CD pipeline с автоматическим запуском SAST
(Bandit) и SCA (pip-audit) на каждый push и pull request. Все проверки
безопасности и автотесты проходят успешно.

\section{Ответы на контрольные вопросы}

\subsection*{1. Почему bcrypt предпочтительнее SHA-256 для хранения паролей?}
SHA-256 — быстрая хэш-функция общего назначения, и эта скорость играет против
защиты: на GPU перебираются миллиарды хэшей в секунду. bcrypt спроектирован
специально для паролей: он намеренно медленный, имеет настраиваемый фактор
сложности (work factor), который можно повышать с ростом мощности железа, и
встроенную случайную соль для каждого пароля. Соль защищает от радужных таблиц и
не даёт одинаковым паролям давать одинаковые хэши.

\subsection*{2. Разница между SAST и DAST.}
SAST (статический анализ) исследует исходный код или байт-код без запуска
приложения и находит потенциально опасные конструкции (небезопасные вызовы,
хардкод секретов). DAST (динамический анализ) тестирует уже запущенное
приложение «снаружи», посылая запросы и наблюдая реакцию, и находит уязвимости
времени выполнения и конфигурации. SAST применяется раньше и видит код, но даёт
ложные срабатывания; DAST видит реальное поведение, но не знает внутренностей.

\subsection*{3. Как работает JWT?}
JWT состоит из трёх частей, разделённых точкой: header (алгоритм и тип), payload
(claims — например \texttt{sub}, \texttt{exp}) и signature. Header и payload —
это base64url от JSON, они не шифруются, а только кодируются. Подпись считается
по header и payload с секретным ключом сервера (для HS256 — HMAC-SHA256). При
проверке сервер пересчитывает подпись по полученным header и payload и сравнивает
с присланной: если ключ верный и данные не менялись — подписи совпадут.
Дополнительно проверяется срок действия \texttt{exp}. Так сервер доверяет токену,
не храня сессию.

\subsection*{4. Чем опасны непроверенные сторонние зависимости?}
В зависимостях (и их транзитивных зависимостях) могут быть известные уязвимости,
которые наследуются приложением. Атакующему достаточно использовать публично
описанную CVE в устаревшей библиотеке. Без регулярного аудита (SCA) такие дыры
остаются незамеченными; отдельные риски — supply-chain атаки и typosquatting
(вредоносные пакеты с похожими именами). Поэтому зависимости фиксируют по версиям
и автоматически сканируют в CI.

\end{document}
